\paragraph{Limitation 2: Workflow-aware caching without a recovery objective.}
Workflow-aware serving already uses such dependencies to improve normal
execution. Parrot exposes inter-request dependencies for application-level
scheduling~\cite{lin2024parrot}. KVFlow guides prefix retention and
prefetching by upcoming agent invocations~\cite{pan2025kvflow}, while Pythia
predicts future workflow steps and prompt contents to manage and stage
KV~\cite{yu2026pythia}. TOPAS jointly selects resident prefixes and admitted
requests by balancing remaining-path progress, reuse, and movement or
preemption costs~\cite{ni2026topas}. These objectives capture normal workflow
progress, but do not explicitly value KV by the WCT increase its retention
prevents after a serving-worker failure.

The gap first concerns \emph{which state is worth retaining}.
Normal reuse depends on future requests, whereas recovery can require
state from an interrupted request even when no later request would reuse
it. For example, KVFlow protects fixed prompts while assigning varying
suffixes the highest eviction priority~\cite{pan2025kvflow}.
Retaining a shared prompt can benefit many future calls, yet leave an
interrupted agent's accumulated task-specific history to be recomputed.
Checkpointing this history can avoid that repeated
work~\cite{strati2024dejavu}. Its protection value therefore depends on
the reconstruction it avoids and the affected branch's slack, in
addition to its expected reuse.

The gap also concerns \emph{where retained state can support recovery}.
CachedAttention and Mooncake already extend KV retention beyond GPU
memory~\cite{gao2024cachedattention,qin2025mooncake}; Pythia writes
future-relevant contexts to shared storage accessible by other
workers~\cite{yu2026pythia}. Such copies can support recovery if they survive
the failure. However, a shared prefix does not eliminate replay of
uncopied per-request history, and KV confined to the failed worker's HBM is
unavailable. Choosing locations for future access alone does not determine
which ongoing histories to protect outside that worker or how quickly
the surviving state must be retrieved.

\label{sec:motivation-long-context}
Figure~\ref{fig:motivation-recompute} shows why this distinction matters
under a storage budget. A surviving copy in a lower-cost tier can already
avoid most reconstruction work. If recovery from that tier fits within
the branch's slack, faster retrieval provides no further WCT benefit
when other execution is unchanged. For a branch with little slack,
the same latency reduction can advance workflow completion.
Meanwhile, normal reuse can justify retaining either branch's KV.
The missing objective is therefore to allocate protection according to
the workflow delay avoided after failure while preserving the value of
KV reuse during normal execution.

